\input{../../../stock/en-tete_v6.tex}

\newcommand{\titrechapitre}{Sujet type Sésame}
\newcommand{\numerochapitre}{1}

\renewcommand{\headrulewidth}{0.4pt}
\renewcommand{\footrulewidth}{0.4pt}
\fancyfoot[L]{Raphaël JONTEF}
\fancyfoot[C]{\thepage}
\fancyfoot[R]{\href{https://www.alveusclub.com/}{Alveus}}
\fancyhead[L]{Mathématiques : Terminale}
\fancyhead[R]{\titrechapitre}

\begin{document}
\input{../../../stock/commands.tex}


\begin{multicols}{2}

\begin{question}{1}{Nombres et dénombrement}
Combien y a-t-il de nombres entiers naturels inférieurs à 200 qui sont à la fois multiples de 4 et de 7 ?
\begin{enumeratebf}
\item 6
\item 7
\item 8
\item 9
\item 10
\end{enumeratebf}
\end{question}

\begin{question}{2}{PPCM et synchronisation}
Deux feux tricolores clignotent respectivement toutes les 48 secondes et toutes les 72 secondes. Après combien de temps clignoteront-ils en même temps pour la première fois ?
\begin{enumeratebf}
\item 120 secondes
\item 144 secondes
\item 168 secondes
\item 192 secondes
\item 240 secondes
\end{enumeratebf}
\end{question}

\begin{question}{3}{Suites arithmétiques}
Dans une suite arithmétique, le premier terme est 5 et le quatrième terme est 17. Quelle est la raison de cette suite ?
\begin{enumeratebf}
\item 2
\item 3
\item 4
\item 5
\item 6
\end{enumeratebf}
\end{question}

\begin{question}{4}{Pourcentages et augmentations}
Un prix augmente de 20\% puis diminue de 10\%. Quel est le pourcentage global d’augmentation par rapport au prix initial ?
\begin{enumeratebf}
\item 8\%
\item 10\%
\item 12\%
\item 14\%
\item 16\%
\end{enumeratebf}
\end{question}

\begin{question}{5}{Problème de vitesse}
Un cycliste roule à 25 km/h pendant 2 heures, puis à 30 km/h pendant 1 heure. Quelle est sa vitesse moyenne sur l’ensemble du trajet ?
\begin{enumeratebf}
\item 26 km/h
\item 27 km/h
\item 28 km/h
\item 29 km/h
\item 30 km/h
\end{enumeratebf}
\end{question}

\begin{question}{6}{Logique de partage}
Un capital de 15 000 € est partagé entre trois associés dans le rapport 3:5:7. Quelle est la part du deuxième associé ?
\begin{enumeratebf}
\item 3 000 €
\item 4 500 €
\item 5 000 €
\item 6 000 €
\item 7 500 €
\end{enumeratebf}
\end{question}

\begin{question}{7}{Nombres premiers}
Quel est le plus petit nombre premier supérieur à 100 ?
\begin{enumeratebf}
\item 101
\item 103
\item 107
\item 109
\item 111
\end{enumeratebf}
\end{question}

\begin{question}{8}{Problème de temps}
Un robinet remplit un bassin en 6 heures, un autre en 9 heures. Combien de temps mettront-ils pour remplir le bassin ensemble ?
\begin{enumeratebf}
\item 3h
\item 3h36
\item 4h
\item 4h30
\item 5h
\end{enumeratebf}
\end{question}

\begin{question}{9}{Logique de codage}
Si A = 1, B = 2, ..., Z = 26, quelle est la valeur du mot « LOGIQUE » ?
\begin{enumeratebf}
\item 78
\item 80
\item 82
\item 84
\item 86
\end{enumeratebf}
\end{question}
\newcolumn

\begin{question}{10}{Problème de mélange}
Un marchand mélange 3 kg de café à 10 €/kg et 2 kg de café à 15 €/kg. Quel est le prix au kg du mélange ?
\begin{enumeratebf}
\item 11 €
\item 11,50 €
\item 12 €
\item 12,50 €
\item 13 €
\end{enumeratebf}
\end{question}

\begin{question}{11}{Nombres et divisibilité}
Quel est le plus petit nombre divisible par 2, 3, 5 et 7 ?
\begin{enumeratebf}
\item 105
\item 140
\item 210
\item 280
\item 315
\end{enumeratebf}
\end{question}

\begin{question}{12}{Problème de distance}
Deux villes sont distantes de 300 km. Un train part de la première ville à 80 km/h, un autre part de la deuxième ville à 70 km/h. Au bout de combien de temps se croiseront-ils ?
\begin{enumeratebf}
\item 1h30
\item 2h
\item 2h30
\item 3h
\item 3h30
\end{enumeratebf}
\end{question}


\begin{question}{13}{Logique de calendrier}
Si le 15 mars est un mardi, quel jour de la semaine sera le 15 avril de la même année ?
\begin{enumeratebf}
\item Mardi
\item Mercredi
\item Jeudi
\item Vendredi
\item Samedi
\end{enumeratebf}
\end{question}

\begin{question}{14}{Problème de proportion}
Dans une recette, il faut 200 g de farine pour 3 œufs. Combien faut-il de farine pour 9 œufs ?
\begin{enumeratebf}
\item 400 g
\item 500 g
\item 600 g
\item 700 g
\item 800 g
\end{enumeratebf}
\end{question}

\begin{question}{15}{Logique de suite}
Trouvez le nombre manquant dans la suite : 2, 6, 12, 20, 30, ?
\begin{enumeratebf}
\item 36
\item 40
\item 42
\item 44
\item 46
\end{enumeratebf}
\end{question}

\begin{question}{16}{Problème de rendement}
Un ouvrier fabrique 120 pièces en 8 heures. Combien de pièces fabriquera-t-il en 5 heures ?
\begin{enumeratebf}
\item 60
\item 70
\item 75
\item 80
\item 90
\end{enumeratebf}
\end{question}

\begin{question}{17}{Logique de fraction}
Quel est le résultat de \( \frac{3}{4} + \frac{5}{6} \) ?
\begin{enumeratebf}
\item \( \frac{8}{10} \)
\item \( \frac{19}{12} \)
\item \( \frac{11}{12} \)
\item \( \frac{17}{12} \)
\item \( \frac{23}{12} \)
\end{enumeratebf}
\end{question}

\begin{question}{18}{Problème de surface}
Un carré a une aire de 144 cm². Quel est son périmètre ?
\begin{enumeratebf}
\item 48 cm
\item 36 cm
\item 24 cm
\item 12 cm
\item 6 cm
\end{enumeratebf}
\end{question}

\begin{question}{19}{Logique de conversion}
Combien y a-t-il de secondes dans 3 heures et 45 minutes ?
\begin{enumeratebf}
\item 12 300
\item 13 500
\item 14 700
\item 15 900
\item 17 100
\end{enumeratebf}
\end{question}

\begin{question}{20}{Problème de moyenne}
Un élève a les notes suivantes : 12, 15, 13, 18. Quelle note doit-il obtenir à son cinquième devoir pour avoir une moyenne de 15 ?
\begin{enumeratebf}
\item 16
\item 17
\item 18
\item 19
\item 20
\end{enumeratebf}
\end{question}

\begin{question}{21}{Logique de volume}
Un cube a une arête de 5 cm. Quel est son volume ?
\begin{enumeratebf}
\item 25 cm³
\item 100 cm³
\item 125 cm³
\item 150 cm³
\item 200 cm³
\end{enumeratebf}
\end{question}

\begin{question}{22}{Problème de pourcentage}
Un article coûte 150 €. Son prix baisse de 20\%. Quel est son nouveau prix ?
\begin{enumeratebf}
\item 100 €
\item 110 €
\item 120 €
\item 130 €
\item 140 €
\end{enumeratebf}
\end{question}

\begin{question}{23}{Logique de temps}
Un train part à 14h30 et arrive à 17h15. Quelle est la durée du trajet ?
\begin{enumeratebf}
\item 2h30
\item 2h45
\item 3h
\item 3h15
\item 3h30
\end{enumeratebf}
\end{question}

\begin{question}{24}{Problème de probabilité}
On lance un dé à 6 faces. Quelle est la probabilité d’obtenir un nombre pair ?
\begin{enumeratebf}
\item 1/6
\item 1/3
\item 1/2
\item 2/3
\item 5/6
\end{enumeratebf}
\end{question}

\begin{question}{25}{Logique de suite complexe}
Trouvez le nombre manquant dans la suite : 1, 4, 9, 16, 25, ?
\begin{enumeratebf}
\item 30
\item 32
\item 34
\item 36
\item 38
\end{enumeratebf}
\end{question}

\end{document}



\end{multicols}
\end{document}